\section[Cote 84 : le produit des algèbres quadratiques, en coordonnées]{Cote 84 : le produit des algèbres quadratiques, en coordonnées\protect\verifie{Folder84}}\label{sec:84}

La cote 84, \emph{[Formes quadratiques] : notes manuscrites (s.d.)}, cent quatre-vingt-douze pages que l'inventaire date \q{[à partir de 1982-vers 1986]}, étudie les formes quadratiques binaires et les algèbres quadratiques au-dessus d'une base quelconque, sans diviser par~$2$. Sa deuxième partie construit le produit $A_1 * A_2$ de deux algèbres quadratiques. Page~47, la page encadre $T_0^2 \equiv \delta \pmod 4$ ; page~48, elle énonce, quand $2$ est régulier, l'équivalence entre revêtements quadratiques et triplets $(L, \delta, T_0)$ (\lot{84}{3}{47}, \lot{84}{3}{48}). Page~63, la page pose \q{sauf erreur} $b = b_1b_2$ et en tire $c$ par $4c = b^2 - \delta$ ; l'encadré écrit \q{$4c_1c_1$}, le membre de droite et le calcul donnent $4c_1c_2$ (\lot{84}{4}{63}). Page~64, elle vérifie par une identité polynomiale que l'élément $u_1b_2 + u_2b_1 + 2u_1u_2$ annule l'équation du produit quand $u_1$ et $u_2$ annulent les leurs (\lot{84}{4}{64}). Page~84, elle calcule des générateurs de $\operatorname{Im}(\sigma_1 \otimes \sigma_2 - \mathrm{id})$ ; page~85, elle énonce la proposition sur les coinvariants, avec un décompte de rangs au-dessus de $X_0 = V(2)$ ; page~86, le corollaire sur les invariants, \q{et inversement si on sait que ceci reste vrai après toute changement de base} (\lot{84}{5}{84}, \lot{84}{5}{85}, \lot{84}{5}{86}). Page~89, le théorème donne $b = b_1b_2$, $c = c_1b_2^2 + c_2b_1^2 - 4c_1c_2$ et $\delta = \delta_1\delta_2$ (\lot{84}{5}{89}) ; page~102, l'image $i(\mathcal U) = 2U_1U_2 + b_1U_2 + b_2U_1$ du générateur dans $A_1 \otimes A_2$ (\lot{84}{6}{102}). Pages~65 à~79, elle introduit les extensions $\chi$-trivialisées et leur produit ; page~78, le cas scindé donne $b = b_1 \otimes b_2$, et une note marginale de la page~79, oblique et en partie illisible, donne le critère d'inversibilité (\lot{84}{4}{78}, \lot{84}{4}{79}). Pages~125 à~131, elle refait le produit en coordonnées pour un quotient $M$ quelconque et commence la relation de cocycle sans l'achever ; page~169, elle met en place les produits poussé et tiré (\lot{84}{7}{125}, \lot{84}{7}{129}, \lot{84}{7}{131}, \lot{84}{9}{169}).

\noindent\emph{Algèbres quadratiques libres.} Pour $b, c \in R$, on note $A(b, c) = R[U]/(U^2 + bU + c)$, de base $(1, U)$, et $\delta(b, c) = b^2 - 4c$. Le changement de variable $V = \varepsilon U + t$, $\varepsilon \in R^\times$, $t \in R$, envoie l'équation de $U$ sur celle de $V$ avec
\[ (b', c') = (\varepsilon b - 2t,\ \varepsilon^2c - \varepsilon bt + t^2). \]
Dans le cas libre, $L = A/R$ est engendré par l'image de $U$, $T_0$ est la réduction de la trace $-b$ modulo~$2$, donc $T_0 \equiv b$ modulo~$2$, et changer de base de $L$ revient à multiplier par $\varepsilon$. Le triplet $(L, \delta, T_0)$ de la page~48 devient, à isomorphisme près, la classe de $(\delta, b \bmod 2)$ sous $(\delta, b) \mapsto (\varepsilon^2\delta, \varepsilon b)$.

\begin{enonce}[Pages 47 à 53, \emph{le théorème des triplets, cas libre}\piste{84-two-regular-triplet-classification}]
Soit $R$ un anneau commutatif.
\begin{enumerate}
\item Les $R$-algèbres $A(b, c)$ et $A(b', c')$ sont isomorphes si et seulement si $(b', c')$ se déduit de $(b, c)$ par un changement de variable $V = \varepsilon U + t$.
\item (Page~47) Pour tous $b, c$, $b^2 - \delta(b, c) = 4c$. Si $\tau^2 - \delta \in 4R$, alors $(\tau + 2a)^2 - \delta \in 4R$ pour tout $a$. Si $\tau^2 - \delta \in 4R$, il existe $c$ avec $\delta(\tau, c) = \delta$.
\item (Page~48) Supposons que $2$ ne divise pas zéro dans $R$. Alors $A(b, c) \simeq A(b', c')$ si et seulement s'il existe $\varepsilon \in R^\times$ avec $\delta(b', c') = \varepsilon^2\delta(b, c)$ et $b' \equiv \varepsilon b$ modulo $2R$. Avec $\varepsilon = 1$ (la base de $L$ fixée), il existe une translation $V = U + t$ de $(b, c)$ sur $(b', c')$ si et seulement si $\delta(b', c') = \delta(b, c)$ et $b' \equiv b$ modulo $2R$, et cette translation est unique.
\item (Page~53, sur $\mathbb Z$) Un entier $\delta$ est de la forme $b^2 - 4c$ si et seulement si $\delta \equiv 0$ ou $1$ modulo~$4$, et $A(b, c) \simeq A(b', c')$ si et seulement si $\delta(b, c) = \delta(b', c')$.
\item Sur $\mathbb F_2$, $(1, 0)$ et $(1, 1)$ ont le même $\delta = 1$ et le même $b$, et $A(1, 0) = \mathbb F_2[U]/(U^2 + U)$ n'est pas isomorphe à $A(1, 1) = \mathbb F_2[U]/(U^2 + U + 1)$.
\end{enumerate}
\end{enonce}

\begin{proof}
(1) Tout élément de $A(b, c)$ s'écrit $a + eU$, et $a + eU = 0$ entraîne $a = e = 0$ : le polynôme unitaire $U^2 + bU + c$ de degré~$2$ ne divise aucun polynôme non nul de degré au plus~$1$, et le reste de la division par lui donne l'écriture. Si $V = \varepsilon U + t$ convient, $\varepsilon U + t$ annule $V^2 + b'V + c'$ dans $A(b, c)$ : le calcul donne $\varepsilon^2(U^2 + bU + c)$. On obtient un morphisme $A(b', c') \to A(b, c)$, $V \mapsto \varepsilon U + t$ ; le changement inverse $U = \varepsilon^{-1}V - \varepsilon^{-1}t$ donne le morphisme réciproque, et les deux composés fixent les générateurs. Réciproquement, soit $\varphi : A(b, c) \to A(b', c')$ un isomorphisme. Écrivons $\varphi^{-1}(V) = t + \varepsilon U$ et $\varphi(U) = s + \eta V$. Alors $U = \varphi^{-1}(\varphi(U)) = (s + \eta t) + \eta\varepsilon U$, donc $\eta\varepsilon = 1$ et $\varepsilon$ est inversible. L'image par $\varphi^{-1}$ de l'équation de $V$, où l'on remplace $U^2$ par $-bU - c$, donne
\[ (t^2 + b't + c' - \varepsilon^2c) + (2t\varepsilon + b'\varepsilon - \varepsilon^2b)\,U = 0, \]
d'où $b' = \varepsilon b - 2t$ après multiplication par $\eta$, puis $c' = \varepsilon^2c - \varepsilon bt + t^2$.

(2) Ce sont des identités : $(\tau + 2a)^2 - \tau^2 = 4(a\tau + a^2)$, et si $\tau^2 - \delta = 4k$, on prend $c = k$.

(3) Un changement de variable donne $\delta(b', c') = \varepsilon^2\delta(b, c)$ en développant, et $b' = \varepsilon b + 2(-t)$ ; cela vaut sur tout anneau. Réciproquement, si $b' = \varepsilon b + 2s$ et $\delta(b', c') = \varepsilon^2\delta(b, c)$, posons $t = -s$. Alors $b' = \varepsilon b - 2t$, et l'égalité des discriminants s'écrit $4c' = 4(\varepsilon^2c - \varepsilon bt + t^2)$. Comme $2$ est régulier, $4$ l'est, et $c' = \varepsilon^2c - \varepsilon bt + t^2$. Le cas $\varepsilon = 1$ est le même calcul ; deux translations $t$, $t'$ vérifient $2t = 2t'$, donc $t = t'$.

(4) Si $b$ est pair, $b^2 - 4c \equiv 0$ ; s'il est impair, $b^2 - 4c \equiv 1$ modulo~$4$. Inversement $\delta = 0^2 - 4(-\delta/4)$ ou $\delta = 1^2 - 4\cdot\frac{1 - \delta}{4}$. Les unités de $\mathbb Z$ sont $\pm 1$, donc $\varepsilon^2 = 1$ et l'énoncé~(3) se réduit à $\delta = \delta'$ et $b' \equiv b$ modulo~$2$. Cette congruence découle de $\delta = \delta'$ : $(b' - b)(b' + b) = 4(c' - c)$ est pair, donc l'un des facteurs l'est, et $b' - b$, $b' + b$ ont la même parité.

(5) Dans $\mathbb F_2$, la seule unité est $1$ et $t^2 = t$ pour tout $t$ ; un changement de variable envoie donc $(1, c)$ sur $(1, c - t + t^2) = (1, c)$. Par~(1), $A(1, 0) \not\simeq A(1, 1)$.
\end{proof}

\begin{remarque}
Ce qui est démontré est le cas libre du théorème de la page~48 : $L$ trivial, l'algèbre donnée par une base $(1, U)$, et la correspondance sur les classes d'isomorphisme. L'équivalence de catégories fibrées, le cas $L$ localement libre non trivial et la reconstruction de $E$, $A$, de la trace, de la norme et du produit par divisions (pages~49 à~52) ne le sont pas. L'énoncé~(1) vaut sur tout anneau et ne demande pas que $2$ soit régulier ; l'hypothèse sert à l'injectivité de $(b, c) \mapsto (\delta, b \bmod 2)$ sur les classes, et l'exemple de $\mathbb F_2$, où $\mathbb F_2 \times \mathbb F_2$ et $\mathbb F_4$ ont les mêmes $\delta$ et $T_0$, montre qu'elle est nécessaire. La surjectivité, (2), vaut sur tout anneau. Sur $\mathbb Z$, on retrouve que les discriminants sont les entiers congrus à $0$ ou $1$ modulo~$4$ et qu'un anneau quadratique est déterminé par son discriminant ; $T_0$ y est déterminé par $\delta$, comme la lecture le dit.
\end{remarque}

\noindent\emph{Conventions.} $R$ est un anneau commutatif quelconque. Les algèbres $A_i$ ont pour bases $(1, U_i)$ avec $U_i^2 + b_iU_i + c_i = 0$, de sorte que $\mathrm{Tr}(U_i) = -b_i$ et que la conjugaison est $\sigma_i(1) = 1$, $\sigma_i(U_i) = -b_i - U_i$ ; on pose $\delta_i = b_i^2 - 4c_i$. C'est la convention de la lecture (pages~98 à~102). Du côté des extensions, $E_i$ a pour base $(e_i, x_i)$, où $e_i$ engendre $L_i$ et $x_i$ est le scindage ; $\psi(e_i) = 0$, $\psi(x_i) = 1$, $T_i = 2x_i - b_ie_i$ et $\sigma_i(x) = \psi(x)T_i - x$, donc $\sigma_i(e_i) = -e_i$ et $\sigma_i(x_i) = x_i - b_ie_i$ (page~84). On écrit $E_1 \otimes E_2 = R^4$ dans la base $(e_1 \otimes e_2, e_1 \otimes x_2, x_1 \otimes e_2, x_1 \otimes x_2)$, $\sigma_1 \otimes \sigma_2$ par le produit de Kronecker des deux matrices, et $E_1 * E_2 = R^2$ dans la base $(e_1 \otimes e_2, x_1 * x_2)$. L'application $E_1 \otimes E_2 \to E_1 * E_2$ est celle de la page~63, $(x_1 + u_1) * (x_2 + u_2) = x_1 * x_2 + u_1b_2 + u_2b_1 + 2u_1u_2$ : elle envoie $(a, a', a'', a''')$ sur $(2a + b_2a' + b_1a'',\ a''')$. Du côté des algèbres, $A = A_1 * A_2$ a pour base $(1, \mathcal U)$, $A_1 \otimes A_2$ la base $(1, U_1, U_2, U_1U_2)$, et $i$ envoie $(\lambda, \mu)$ sur $(\lambda, \mu b_2, \mu b_1, 2\mu)$.

\begin{enonce}[Pages 63, 64, 89 et 102, \emph{la formule du produit}\piste{84-product-quadratic-algebras-formula}]
Soient $A_1$, $A_2$ deux $R$-algèbres commutatives et $u_i \in A_i$ avec $u_i^2 + b_iu_i + c_i = 0$, $b_i, c_i \in R$. Posons $U_1 = u_1 \otimes 1$, $U_2 = 1 \otimes u_2$ dans $A_1 \otimes_R A_2$, $b = b_1b_2$ et $c = c_1b_2^2 + c_2b_1^2 - 4c_1c_2$.
\begin{enumerate}
\item (Page~64) Dans tout anneau commutatif, avec $Q_i = u_i^2 + b_iu_i + c_i$ et $u = 2u_1u_2 + b_1u_2 + b_2u_1$,
\[ u^2 + bu + c = b_2^2Q_1 + b_1^2Q_2 + 4\bigl(Q_1Q_2 - c_1Q_2 - c_2Q_1\bigr). \]
\item (Page~102) L'élément $\mathcal U = 2U_1U_2 + b_1U_2 + b_2U_1$ vérifie $\mathcal U^2 + b\,\mathcal U + c = 0$ dans $A_1 \otimes_R A_2$ ; il existe donc un unique morphisme de $R$-algèbres $R[\mathcal U]/(\mathcal U^2 + b\,\mathcal U + c) \to A_1 \otimes_R A_2$ envoyant $\mathcal U$ sur cet élément.
\item (Pages~63 et~89) $b^2 - 4c = \delta_1\delta_2$, et $c_1\delta_2 + c_2\delta_1 + 4c_1c_2 = c_1b_2^2 + c_2b_1^2 - 4c_1c_2$.
\end{enumerate}
\end{enonce}

\begin{proof}
(1) et (3) sont des identités entre polynômes à coefficients entiers ; on les vérifie en développant. (2) Les inclusions $A_1 \to A_1 \otimes A_2$ et $A_2 \to A_1 \otimes A_2$ sont des morphismes de $R$-algèbres ; elles transportent les équations de $u_1$ et $u_2$ en $U_1^2 + b_1U_1 + c_1 = 0$ et $U_2^2 + b_2U_2 + c_2 = 0$. L'identité~(1), appliquée dans l'anneau $A_1 \otimes A_2$ à $U_1$ et $U_2$, donne $\mathcal U^2 + b\,\mathcal U + c = 0$, puisque $Q_1 = Q_2 = 0$. La propriété universelle de $R[\mathcal U]/(\mathcal U^2 + b\,\mathcal U + c)$ donne le morphisme.
\end{proof}

\begin{remarque}
L'identité de la page~64 vaut sans hypothèse sur $u_1$, $u_2$, dans tout anneau commutatif, et la vérification ne divise jamais par~$2$ : le produit est défini sur $\mathbb Z$. La lecture de \q{$4c_1c_1$} en $4c_1c_2$ est confirmée : avec $4c_1c_2$, l'identité~(1) et la relation $b^2 - 4c = \delta_1\delta_2$ tiennent, et la première forme de l'encadré est égale à la seconde. Ce qui est démontré est la partie en coordonnées de la piste ; la construction intrinsèque de $X_1 * X_2$ et l'unicité de la forme quadratique de la page~89 ne le sont pas.
\end{remarque}

\begin{enonce}[Pages 84 et 85, \emph{proposition}\piste{84-product-coinvariants-one-etale}]
Dans les coordonnées ci-dessus, l'application induite
\[ (E_1 \otimes E_2)/\operatorname{Im}(\sigma_1 \otimes \sigma_2 - \mathrm{id}) \to E_1 * E_2 \]
est bien définie, et elle est bijective si et seulement si $2R + b_1R + b_2R = R$.
\end{enonce}

\begin{proof}
On calcule $\sigma_1 \otimes \sigma_2 - \mathrm{id}$ sur la base : $e_1 \otimes e_2 \mapsto 0$, $e_1 \otimes x_2 \mapsto (b_2, -2, 0, 0)$, $x_1 \otimes e_2 \mapsto (b_1, 0, -2, 0)$, $x_1 \otimes x_2 \mapsto (b_1b_2, -b_1, -b_2, 0)$ ; ce sont les générateurs de la page~84. L'application $(a, a', a'', a''') \mapsto (2a + b_2a' + b_1a'', a''')$ les annule, d'où l'application induite.

Elle a même image que $E_1 \otimes E_2 \to E_1 * E_2$. Si elle est surjective, $(1, 0)$ est atteint, donc $2a + b_2a' + b_1a'' = 1$ pour des $a, a', a'' \in R$, et $2R + b_1R + b_2R = R$. Réciproquement, si $2g_0 + b_1g_1 + b_2g_2 = 1$, l'image de $(y_0g_0, y_0g_2, y_0g_1, y_1)$ est $(y_0, y_1)$.

Sous la condition, montrons que le noyau est dans $\operatorname{Im}(\sigma_1 \otimes \sigma_2 - \mathrm{id})$ ; l'injectivité de l'application induite en résulte. Écrivons $f = (2, b_2, b_1)$ et $g = (g_0, g_2, g_1)$, de sorte que $f \cdot g = 1$. Un élément du noyau est $(a_0, a_1, a_2, 0)$ avec $f \cdot a = 0$. Les relations de Koszul $\kappa_{jk} = f_k\varepsilon_j - f_j\varepsilon_k$ de $f$ dans $R^3$ sont $\kappa_{01} = (b_2, -2, 0)$, $\kappa_{02} = (b_1, 0, -2)$ et $\kappa_{12} = (0, b_1, -b_2)$ ; les deux premières sont les images de $e_1 \otimes x_2$ et $x_1 \otimes e_2$, la troisième celle de $x_1 \otimes x_2 - b_1\, e_1 \otimes x_2$. Or
\[ a = \sum_{j < k} (g_ka_j - g_ja_k)\,\kappa_{jk}, \]
car la $i$-ème coordonnée du second membre est $a_i (f \cdot g) - g_i (f \cdot a) = a_i$. L'élément $a$ est donc l'image d'une combinaison explicite des vecteurs de base.
\end{proof}

\begin{enonce}[Page 86, \emph{corollaire}\piste{84-product-coinvariants-one-etale}]
Dans les coordonnées ci-dessus, disons que $A \to (A_1 \otimes A_2)^{\sigma_1 \otimes \sigma_2}$ est un isomorphisme si $i$ est injectif et si son image est le sous-module des invariants de $\sigma_1 \otimes \sigma_2$.
\begin{enumerate}
\item Si $2R + b_1R + b_2R = R$, c'est un isomorphisme.
\item C'est un isomorphisme après tout changement de base $R \to S$ si et seulement si $2R + b_1R + b_2R = R$.
\item Sur $R = \mathbb Z[X]$, avec $b_1 = X$ et $b_2 = 0$, c'est un isomorphisme, alors que $2$, $X$ et $0$ engendrent un idéal propre.
\end{enumerate}
\end{enonce}

\begin{proof}
Un vecteur $(p, q, r, s)$ (coefficients de $1$, $U_1$, $U_2$, $U_1U_2$) est invariant si et seulement si $-b_1q - b_2r + b_1b_2s = 0$, $2q = b_2s$ et $2r = b_1s$. L'image $(\lambda, \mu b_2, \mu b_1, 2\mu)$ de $i$ vérifie ces relations.

(1) Soit $2g_0 + b_1g_1 + b_2g_2 = 1$. Si $i(\lambda, \mu) = 0$, alors $\lambda = 0$ et $\mu = \mu(2g_0 + b_1g_1 + b_2g_2) = 0$. Soit $(p, q, r, s)$ invariant. Des relations on tire $b_1q = b_2r$, car $b_1b_2s = 2b_1q$ (on multiplie $2q = b_2s$ par $b_1$). Posons $\mu = sg_0 + rg_1 + qg_2$. Alors $2\mu = s$, $b_1\mu = r$ et $b_2\mu = q$ ; par exemple $b_2\mu - q = g_0(b_2s - 2q) + g_1(b_2r - b_1q) = 0$, en remplaçant $q$ par $q(2g_0 + b_1g_1 + b_2g_2)$. Donc $(p, q, r, s) = i(p, \mu)$.

(2) La condition passe à tout $S$ : l'image de $1 = 2g_0 + b_1g_1 + b_2g_2$ est une relation du même type, et (1) s'applique. Réciproquement, si l'idéal $I = (2, b_1, b_2)$ n'est pas $R$, prenons $S = R/I$, qui n'est pas nul. Dans $S$, $2 = b_1 = b_2 = 0$, et le vecteur $U_2 = (0, 0, 1, 0)$ est invariant. S'il était $i(\lambda, \mu)$, sa coordonnée sur $U_2$ serait $\mu b_1 = 0$, et l'on aurait $1 = 0$ dans $S$.

(3) Avec $b_1 = X$, $b_2 = 0$, les relations deviennent $2q = 0$, $2r = Xs$. L'anneau $\mathbb Z[X]$ est intègre, donc $q = 0$ et $i$ est injectif ($2\mu = 0$ entraîne $\mu = 0$). L'élément $2$ est premier dans $\mathbb Z[X]$ et ne divise pas $X$ ; de $2 \mid Xs$ on tire $s = 2\mu$, puis $r = X\mu$ en simplifiant par~$2$, et $(p, q, r, s) = i(p, \mu)$. Enfin, si $2g_0 + Xg_1 = 1$, le coefficient constant donne $2\,g_0(0) = 1$ dans $\mathbb Z$, ce qui est impossible.
\end{proof}

\begin{remarque}
La proposition de la page~85 tient en coordonnées sur tout anneau commutatif, dans les deux sens, sans changement de base, sans hypothèse de rang et sans la réduction au cas $2 = 0$ que la page annonce. Le noyau est engendré par les relations de Koszul de la suite unimodulaire $(2, b_2, b_1)$. Pour le corollaire de la page~86, la condition suffit sur $R$ lui-même, et la réciproque demande l'hypothèse de la page, \q{après toute changement de base} : l'exemple de $\mathbb Z[X]$ montre que l'isomorphisme sur $R$ seul n'entraîne pas la condition. Un changement de base vers $R/(2, b_1, b_2)$ suffit. La formulation ponctuelle (en tout point de $V(2)$, l'un des $T_{i,0}$ engendre $L_{i,0}$) n'est pas démontrée ; elle équivaut à la condition en coordonnées dès que $L_1$, $L_2$ sont triviaux et munis de scindages, ce que la page dit aussi. Le passage de l'énoncé sur les coinvariants à l'énoncé sur les invariants par dualité n'est pas formalisé : les deux sont démontrés séparément, en coordonnées.
\end{remarque}

\noindent\emph{Extensions $\chi$-scindées en coordonnées.} Soit $\chi \in R$. On prend une extension scindée $0 \to L \to M \oplus L \to M \to 0$, $\alpha(\xi) = (0, \xi)$, $\psi(\eta, \xi) = \eta$, de $R$-modules quelconques. Un $\chi$-scindage est un couple $(\pi, \varpi)$, $\pi : M \oplus L \to L$, $\varpi : M \to M \oplus L$, avec $\pi\alpha = \chi$, $\psi\varpi = \chi$ et $\alpha\pi + \varpi\psi = \chi$ (page~169). Pour $b : M \to L$, on pose $\pi_b(\eta, \xi) = b(\eta) + \chi\xi$ et $\varpi_b(\eta) = (\chi\eta, -b(\eta))$. La transvection $T_w$, pour $w : M \to L$, est $(\eta, \xi) \mapsto (\eta, \xi + w(\eta))$. Le produit de deux objets $b_1 : M_1 \to L_1$ et $b_2 : M_2 \to L_2$ est $b_1 \otimes b_2 : M_1 \otimes M_2 \to L_1 \otimes L_2$ (page~78, \q{Il n'y a pas plus simple !}, et page~129). Pour $M = L = R$, un objet est un scalaire $b$, et l'objet unité est $b = 1$, $\pi(u, \lambda) = \lambda + \chi u$ (page~78).

\begin{enonce}[Pages 77--79 et 123--131, extensions $\chi$-scindées\piste{84-chi-trivialised-extension-product}]
\begin{enumerate}
\item Le couple $(\pi_b, \varpi_b)$ est un $\chi$-scindage. Réciproquement, tout $\chi$-scindage $(\pi, \varpi)$ de $M \oplus L$ est de cette forme, avec $b = \pi|_M$.
\item Il existe un endomorphisme $f$ de $M \oplus L$ induisant l'identité sur $L$ et sur $M$ et tel que $\pi_{b'} \circ f = \pi_b$ si et seulement si $b = b' + \chi w$ pour un $w : M \to L$ ; c'est alors $T_w$. On a aussi $T_w \circ \varpi_b = \varpi_{b'}$ et $T_w \circ T_{w'} = T_{w + w'}$.
\item (Page~129) Pour $u_i : M_i \to L_i$, $(b_1 + \chi u_1) \otimes (b_2 + \chi u_2) = b_1 \otimes b_2 + \chi u$ avec $u = b_1 \otimes u_2 + u_1 \otimes b_2 + \chi\, u_1 \otimes u_2$.
\item (Page~131) Si l'on change d'abord par $u_i$, puis par $u'_i$ en partant de $b'_i = b_i + \chi u_i$, la somme des deux données est celle du changement par $u_i + u'_i$ : $u'' = u + u'$.
\item (Page~78) Le produit est associatif, commutatif et unitaire : à travers les isomorphismes d'associativité et de commutativité des produits tensoriels, $(b_1 \otimes b_2) \otimes b_3$ correspond à $b_1 \otimes (b_2 \otimes b_3)$ et $b_1 \otimes b_2$ à $b_2 \otimes b_1$ ; à travers $R \otimes M \simeq M$ et $R \otimes L \simeq L$, $\mathrm{id}_R \otimes b$ correspond à $b$.
\item (Page~79) Pour $M = L = R$, disons $b \sim b'$ s'il existe $\varepsilon \in R^\times$ et $w \in R$ avec $\varepsilon b = b' + \chi w$ (un isomorphisme qui agit par $\varepsilon$ sur $L$). La relation est compatible au produit, et la classe de $b$ est inversible, c'est-à-dire $bb' \sim 1$ pour un $b'$, si et seulement si $bR + \chi R = R$, si et seulement si $b$ est inversible dans $R/\chi R$ : la section $T_0 = -b_0$ engendre $L_0$ sur $V(\chi)$.
\end{enumerate}
\end{enonce}

\begin{proof}
(1) Les trois relations se vérifient sur $(\eta, 0)$ et $(0, \xi)$ ; par exemple $\alpha\pi_b(\eta, \xi) + \varpi_b(\eta) = (0, b\eta + \chi\xi) + (\chi\eta, -b\eta) = \chi(\eta, \xi)$. Si $(\pi, \varpi)$ est un $\chi$-scindage, $\pi(0, \xi) = \chi\xi$ donne $\pi = \pi_b$. La relation $\alpha\pi + \varpi\psi = \chi$ appliquée à $(\eta, 0)$ donne $\varpi(\eta) = \chi(\eta, 0) - (0, b\eta)$, et $\varpi = \varpi_b$.

(2) Si $f$ induit l'identité sur $L$ et sur $M$, alors $f(0, \xi) = (0, \xi)$ et la première coordonnée de $f(\eta, 0)$ est $\eta$, donc $f = T_w$ avec $w(\eta)$ la seconde coordonnée de $f(\eta, 0)$. On a $\pi_{b'}(T_w(\eta, \xi)) = b'\eta + \chi\xi + \chi w\eta = \pi_{b' + \chi w}(\eta, \xi)$, d'où la condition $b = b' + \chi w$, en évaluant en $(\eta, 0)$. Inversement $T_w$ convient. Le calcul de $T_w \circ \varpi_b$ et de $T_w \circ T_{w'}$ est direct.

(3) et (4) Le produit tensoriel d'applications linéaires est bilinéaire ; on développe. Pour (4), on compare $u + u'$, où $u' = (b_1 + \chi u_1) \otimes u'_2 + u'_1 \otimes (b_2 + \chi u_2) + \chi\, u'_1 \otimes u'_2$, avec $(b_1 \otimes (u_2 + u'_2)) + ((u_1 + u'_1) \otimes b_2) + \chi\,(u_1 + u'_1) \otimes (u_2 + u'_2)$ ; les deux développements ont les mêmes neuf termes.

(5) Ce sont la naturalité des isomorphismes d'associativité et de commutativité, qui sont dans mathlib, et un calcul sur les tenseurs purs pour l'unité.

(6) Si $\varepsilon_i b_i = b'_i + \chi w_i$, alors $\varepsilon_1\varepsilon_2 b_1b_2 = b'_1b'_2 + \chi(b'_1w_2 + w_1b'_2 + \chi w_1w_2)$. Si $\varepsilon bb' = 1 + \chi w$, alors $b(\varepsilon b') + \chi(-w) = 1$. Si $bg + \chi h = 1$, alors $bg = 1 + \chi(-h)$, et $bg \sim 1$ avec $\varepsilon = 1$. Enfin $bR + \chi R = R$ équivaut à l'existence de $g$ avec $bg \equiv 1$ modulo $\chi R$.
\end{proof}

\begin{remarque}
Ce qui est démontré est la partie en coordonnées de la piste, pour une extension scindée de modules quelconques, sans diviser par $\chi$ ni par~$2$. Le recollement de ces coordonnées en un produit $E_1 * E_2$ sur une base quelconque, la propriété universelle des applications bi-affines (pages~75 et~76) et la compatibilité des contraintes entre elles (la quatrième compatibilité de la page~177 n'est pas écrite) ne le sont pas. La relation de cocycle de la page~131, que la page commence sans l'achever, est vraie avec les signes $b'_i = b_i + \chi u_i$. Le critère d'inversibilité de la note marginale de la page~79 est démontré dans le cas $L = M = R$ : l'inverse de la classe de $b$ est la classe de $g$ dès que $bg \equiv 1$ modulo~$\chi$, et l'on n'a besoin d'aucune unité $\varepsilon$ autre que $1$ pour le construire.
\end{remarque}

\begin{enonce}[Pages 125, 127 et 169, \emph{produit poussé et produit tiré}\piste{84-chi-trivialised-extension-product}]
En coordonnées, le produit poussé porte le $\chi$-scindage de $b = b_1 \otimes b_2$ et le produit tiré celui de $-b$ (page~125).
\begin{enumerate}
\item Si $\gamma\chi = 2$, la transvection $T_{\gamma b}$ envoie le $\chi$-scindage $\pi_b$ sur $\pi_{-b}$ : $\pi_{-b} \circ T_{\gamma b} = \pi_b$.
\item (Page~125, \q{commute !}) Si $\gamma\chi = 2$, pour un changement de scindage qui remplace $b$ par $b + \chi u$ du côté poussé et $-b$ par $-b - \chi u$ du côté tiré, $T_{\gamma b} \circ T_u = T_{-u} \circ T_{\gamma(b + \chi u)}$.
\item Pour les objets unités ($M = L = R$, $b_1 = b_2 = \mathrm{id}$), il existe un endomorphisme de $R \oplus R$ qui induit l'identité sur $L$ et sur $M$ et envoie le scindage poussé sur le scindage tiré si et seulement s'il existe $\gamma$ avec $\gamma\chi = 2$. Sur $\mathbb Z$, avec $\chi = 4$, il n'en existe pas.
\end{enumerate}
\end{enonce}

\begin{proof}
(1) Par l'énoncé précédent, $\pi_{-b} \circ T_{\gamma b} = \pi_{-b + \chi\gamma b} = \pi_{-b + 2b} = \pi_b$. (2) Les transvections se composent par addition : il faut $\gamma b + u = -u + \gamma b + \gamma\chi u$, ce qui est $\gamma\chi u = 2u$. (3) Par l'énoncé précédent, un tel endomorphisme existe si et seulement si $\mathrm{id} = -\mathrm{id} + \chi w$ pour un $w : R \to R$ ; en $1$, cela donne $w(1)\chi = 2$, et inversement $w = \gamma\,\mathrm{id}$ convient. Sur $\mathbb Z$, $4\gamma = 2$ n'a pas de solution.
\end{proof}

\begin{remarque}
L'hypothèse $\gamma\chi = 2$, que la lecture prend à la page~127 pour l'isomorphisme de la page~169, est suffisante, et elle est nécessaire dès les objets unités, parmi les isomorphismes qui induisent l'identité sur $L_1 \otimes L_2$ et $M_1 \otimes M_2$ : la comparaison demande que $2$ soit multiple de $\chi$. Le signe $-b$ du produit tiré est pris de la page~125 ; la construction du produit tiré comme produit fibré (page~169) et l'identification de la flèche du milieu du diagramme de la page~167 avec l'opérateur $F_{11}$ ne sont pas formalisées.
\end{remarque}

Ne sont pas démontrés ici : l'équivalence de catégories fibrées du théorème de la page~48 et son cas $L$ non trivial, la reconstruction de $E$, $A$, de la trace, de la norme et du produit par divisions (pages~49 à~52), la construction intrinsèque du produit $X_1 * X_2$ et des extensions $\chi$-scindées sur une base quelconque (pages~65 à~79), la propriété universelle des applications bi-affines (pages~75 et~76), le critère d'inversibilité pour $L$ inversible non trivial, la construction du produit tiré comme produit fibré et l'opérateur $F_{11}$ (pages~141 à~169), l'unicité de la forme quadratique du théorème de la page~89, la formulation ponctuelle de la proposition de la page~85 et du corollaire de la page~86, l'énoncé du corollaire pour $\underline{\mathrm{Hom}}(-, M)$, la trace et la norme de $A_1 \otimes A_2$ sur $A$ (pages~104 et~105), la quatrième partie du dossier et le reste de la troisième.
